% Part: incompleteness
% Chapter: representability-in-q
% Section: minimization-representable

\documentclass[../../../include/open-logic-section]{subfiles}

\begin{document}

\olfileid[pt]{inc}{req}{min}
\olsection{A minimização regular é representável em $\Th{Q}$}

Consideremos a busca ilimitada. Suponha que $g(x, z)$ seja regular e
representável em $\Th{Q}$, digamos, pela !!{formula}~$!A_g(x, z, y)$.
Seja $f$ definida por $f(z) = \umin{x}{[g(x, z) = 0]}$. Gostaríamos de
encontrar !!a{formula}~$!A_f(z, y)$ que represente~$f$. O valor
de~$f(z)$ é o número~$x$ que (a) satisfaz $g(x, z) = 0$ e (b) é o menor
com essa propriedade, isto é, para todo $w < x$, $g(w, z) \neq 0$.
Assim, a escolha a seguir é natural:
\[
!A_f(z,y) \ident !A_g(y, z, \Obj 0) \land \lforall[w][(w < y \lif \lnot
  !A_g(w, z, \Obj 0))].
\]
No caso geral, naturalmente, teríamos de substituir $z$ por
$z_0$, \dots, $z_k$.

Mais uma vez, a prova envolverá alguns lemas sobre coisas que $\Th{Q}$
é forte o suficiente para provar.

\begin{lem}
\ollabel{lem:succ} Para toda !!{constant}~$a$ e todo número
natural~$n$,
\[
\Th{Q} \Proves \eq[(a' + \num n)][(a + \num n)'].
\]
\end{lem}

\begin{proof}
A prova, como de costume, é por indução sobre $n$. No caso base,
$n = 0$, precisamos mostrar que $\Th{Q}$ prova
$\eq[(a' + \Obj 0)][(a + \Obj 0)']$. Mas temos:
\begin{align}
  \Th{Q} & \Proves \eq[(a' + \Obj 0)][a'] \quad \text{pelo axioma $Q_4$}
  \ollabel{step1}\\
  \Th{Q} & \Proves  \eq[(a + \Obj 0)][a] \quad \text{pelo axioma $Q_4$}
  \ollabel{step2} \\
  \Th{Q} & \Proves \eq[(a + \Obj 0)'][a'] \quad \text{por \olref{step2}}
  \ollabel{step3} \\
  \Th{Q} & \Proves \eq[(a' + \Obj 0)][(a + \Obj 0)'] \quad
  \text{por \olref{step1} e \olref{step3}}\notag
\end{align}
No passo de indução, podemos supor que mostramos que $\Th{Q}
\Proves \eq[(a' + \num n)][(a + \num n)']$. Como $\num{n+1}$ é
$\num{n}'$, precisamos mostrar que $\Th{Q}$ prova
$\eq[(a' + \num{n}')][(a + \num{n}')']$. Temos:
\begin{align}
  \Th{Q} & \Proves \eq[(a' + \num n')][(a' + \num n)'] \quad
  \text{pelo axioma $!Q_5$} \ollabel{step5}\\
  \Th{Q} & \Proves \eq[(a' + \num n')][(a + \num n')'] \quad
  \text{hipótese de indução} \ollabel{step6}\\
  \Th{Q} & \Proves \eq[(a' + \num n)'][(a + \num n')'] \quad
  \text{por \olref{step5} e \olref{step6}.} \notag
\end{align}
\end{proof}

Vale novamente mencionar que isso é mais fraco do que dizer que
$\Th{Q}$ prova $\lforall[x][\lforall[y][(x' + y) = (x + y)']]$.
Embora essa !!{sentence} seja verdadeira em~$\Struct{N}$, $\Th{Q}$ não
a prova.

\begin{lem}
\ollabel{lem:less-zero}
$\Th{Q} \Proves \lforall[x][\lnot x < \Obj 0]$.
\end{lem}

\begin{proof}
Apresentamos a prova informalmente (isto é, dando apenas indicações de
como construir a !!{derivation} formal).

Precisamos provar $\lnot a < \Obj 0$ para um~$a$ arbitrário. Pela
definição de $<$, precisamos provar
$\lnot \lexists[y][\eq[(y'+a)][\Obj 0]]$ em $\Th{Q}$. Suporemos
$\lexists[y][\eq[(y'+a)][\Obj 0]]$ e provaremos uma contradição.
Suponha $\eq[(b'+a)][\Obj 0]$. Usando $!Q_3$, temos
$\eq[a][\Obj 0] \lor \lexists[y][\eq[a][y']]$. Distinguimos casos.

Caso 1: $\eq[a][\Obj 0]$ vale. De $\eq[(b'+a)][\Obj 0]$, obtemos
$\eq[(b' + \Obj 0)][\Obj 0]$. Pelo axioma~$!Q_4$ de $\Th{Q}$, temos
$\eq[(b' + \Obj 0)][b']$ e, portanto, $\eq[b'][\Obj 0]$. Mas, pelo
axioma~$!Q_2$, também temos $\eq/[b'][\Obj 0]$, uma contradição.

Caso 2: para algum $c$, $\eq[a][c']$. Nesse caso, temos
$\eq[(b' + c')][\Obj 0]$. Pelo axioma~$!Q_5$, temos
$\eq[(b' + c)'][\Obj 0]$, contradizendo novamente o axioma~$Q_2$.
\end{proof}

\begin{lem}
\ollabel{lem:less-nsucc}
Para todo número natural~$n$,
  \[
  \Th{Q} \Proves
  \lforall[x][(x < \num {n+1} \lif (\eq[x][\Obj 0] \lor \dots \lor
    \eq[x][\num n]))].
  \]
\end{lem}

\begin{proof}
Usamos indução sobre~$n$. Consideremos o caso base, em que $n = 0$.
Nesse caso, precisamos mostrar $a < \num 1 \lif \eq[a][\Obj 0]$ para
um~$a$ arbitrário. Suponha $a < \num 1$. Então, pelo axioma definidor
de $<$, temos $\lexists[y][\eq[(y'+a)][\Obj 0']]$ (pois
$\num 1 \ident \Obj 0'$).

Suponha que $b$ tenha essa propriedade, isto é, que tenhamos
$\eq[(b'+a)][\Obj 0']$. Precisamos mostrar $\eq[a][\Obj 0]$. Pelo
axioma~$!Q_3$, temos ou $\eq[a][\Obj 0]$, ou existe um $c$ tal que
$\eq[a][c']$. No primeiro caso, não há nada a mostrar. Suponha, então,
$\eq[a][c']$. Temos $\eq[(b' + c')][\Obj 0']$. Pelo axioma~$!Q_5$ de
$\Th{Q}$, temos $\eq[(b'+c)'][\Obj 0']$. Pelo axioma~$!Q_1$, temos
$\eq[(b' + c)][\Obj 0]$. Mas isso significa, pelo axioma~$!Q_8$, que
$c < \Obj 0$, contradizendo \olref{lem:less-zero}.

Passemos ao passo de indução. Provamos o caso para $n+1$, supondo o
caso para~$n$. Suponha, então, $a < \num {n+2}$. Usando novamente
$!Q_3$, podemos distinguir dois casos: $\eq[a][\Obj 0]$ e, para algum
$b$, $\eq[a][b']$. No primeiro, $\eq[a][\Obj 0] \lor \dots \lor
\eq[a][\num{n+1}]$ decorre trivialmente. No segundo, temos
$b' < \num {n+2}$, isto é, $b' < \num{n+1}'$. Pelo axioma~$!Q_8$,
para algum $c$, $\eq[(c'+b')][\num{n+1}']$. Pelo axioma $!Q_5$,
$\eq[(c'+b)'][\num{n+1}']$. Pelo axioma~$!Q_1$,
$\eq[(c'+b)][\num{n+1}]$ e, portanto, $b < \num{n+1}$ pelo
axioma~$!Q_8$. Pela hipótese de indução,
$\eq[b][\Obj 0] \lor \dots \lor \eq[b][\num{n}]$. Daí obtemos
$\eq[b'][\Obj 0'] \lor \dots \lor \eq[b'][\num{n}']$ pela lógica e,
portanto, $\eq[a][\num{1}] \lor \dots \lor \eq[a][\num{n+1}]$, pois
$\eq[a][b']$.
\end{proof}

\begin{lem}
  \ollabel{lem:trichotomy} Para todo número natural~$m$,
  \[
  \Th{Q} \Proves
  \lforall[y][((y < \num{m} \lor \num{m} < y) \lor \eq[y][\num{m}])].
  \]
\end{lem}

\begin{proof}
Por indução sobre~$m$. Primeiro, considere o caso $m=0$. $\Th{Q} \Proves
\lforall[y][(\eq[y][\Obj 0] \lor \lexists[z][\eq[y][z']])]$ por~$!Q_3$.
Seja $a$ arbitrário. Então ou $\eq[a][\Obj 0]$, ou, para algum~$b$,
$\eq[a][b']$. No primeiro caso, também temos
$(a < \Obj 0 \lor \Obj 0 < a) \lor \eq[a][\Obj 0]$. Mas, se
$\eq[a][b']$, então $\eq[(b' + \Obj 0)][(a + \Obj 0)]$ pela lógica
de~$\eq$. Por $!Q_4$, $\eq[(a + \Obj 0)][a]$; logo, temos
$\eq[(b' + \Obj 0)][a]$ e, portanto,
$\lexists[z][\eq[(z' + \Obj 0)][a]]$. Pela definição de $<$ em $!Q_8$,
$\Obj 0 < a$. Se $\Obj 0 < a$, então também
$(\Obj 0 < a \lor a < \Obj 0) \lor \eq[a][\Obj 0]$.

Suponha agora que tenhamos
\begin{align*}
  \Th{Q} & \Proves \lforall[y][((y < \num{m} \lor \num{m} < y) \lor
    \eq[y][\num{m}])]
  \intertext{e queremos mostrar}
  \Th{Q} & \Proves \lforall[y][((y < \num{m+1} \lor \num{m+1} < y) \lor
    \eq[y][\num{m+1}])]
\end{align*}
Seja $a$ arbitrário. Por $!Q_3$, ou $\eq[a][\Obj 0]$, ou, para
algum~$b$, $\eq[a][b']$. No primeiro caso, temos
$\eq[\num{m}' + a][\num{m+1}]$ por $!Q_4$ e, portanto,
$a < \num{m+1}$ por $!Q_8$.

Considere agora o segundo caso, $\eq[a][b']$. Pela hipótese de indução,
$(b < \num{m} \lor \num{m} < b) \lor \eq[b][\num{m}]$.

A primeira disjunção, $b < \num{m}$, é equivalente (por $!Q_8$) a
$\lexists[z][\eq[(z' + b)][\num{m}]]$. Suponha que $c$ tenha essa
propriedade. Se $\eq[(c' + b)][\num{m}]$, então também
$\eq[(c' + b)'][\num{m}']$. Por $!Q_5$,
$\eq[(c' + b)'][(c' + b')]$. Portanto,
$\eq[(c' + b')][\num{m}']$. Obtemos
$\lexists[u][\eq[(u' + b')][\num{m+1}]]$ generalizando
existencialmente sobre~$c'$ e lembrando que
$\num{m}' \ident \num{m+1}$. Logo, se $b < \num{m}$, então
$b' < \num{m+1}$ e, portanto, $a < \num{m+1}$.

Suponha agora $\num{m} < b$, isto é,
$\lexists[z][\eq[(z' + \num{m})][b]]$. Suponha que $c$ seja tal~$z$,
isto é, $\eq[(c' + \num{m})][b]$. Pela lógica,
$\eq[(c' + \num{m})'][b']$. Por $!Q_5$,
$\eq[(c' + \num{m}')][b']$. Como $\eq[a][b']$ e
$\num{m}' \ident \num{m+1}$, $\eq[(c' + \num{m+1})][a]$. Por $!Q_8$,
$\num{m+1} < a$.

Por fim, suponha $\eq[b][\num{m}]$. Então, pela lógica,
$\eq[b'][\num{m}']$ e, portanto, $\eq[a][\num{m+1}]$.

Assim, de cada disjunção do caso para~$m$ e~$b$, podemos obter a
disjunção correspondente para~$m+1$ e~$a$.
\end{proof}
  
\begin{prop}
\ollabel{prop:rep-minimization}
Se $!A_g(x, z, y)$ representa $g(x, z)$ em~$\Th{Q}$, então
\[
!A_f(z,y) \ident !A_g(y, z, \Obj 0) \land \lforall[w][(w < y \lif \lnot
  !A_g(w, z, \Obj 0))]
\]
representa $f(z) = \umin{x}{[g(x, z) = 0]}$.
\end{prop}

\begin{proof}
Primeiro mostramos que, se $f(n) = m$, então
$\Th{Q} \Proves !A_f(\num n, \num m)$, isto é,
\begin{align}
  \Th{Q} & \Proves !A_g(\num{m}, \num{n}, \Obj 0) \land \lforall[w][(w <
    \num{m} \lif \lnot !A_g(w, \num{n}, \Obj 0))]. \notag
  \intertext{Como $!A_g(x, z, y)$ representa $g(x, z)$ e $g(m, n) =
    0$ se $f(n) = m$, temos}
\Th{Q} & \Proves !A_g(\num{m}, \num{n}, \Obj 0). \notag
\intertext{Se $f(n) = m$, então, para todo $k < m$, $g(k, n) \neq 0$. Logo,}
\Th{Q} & \Proves \lnot !A_g(\num{k}, \num{n}, \Obj 0). \notag
\intertext{Obtemos}
\Th{Q} & \Proves \lforall[w][(w < \num{m} \lif \lnot
  !A_g(w, \num{n}, \Obj 0))]. \ollabel{rep-less}
\end{align}
por \olref{lem:less-zero} no caso $m = 0$ e por
\olref{lem:less-nsucc} caso contrário.

Mostremos agora que, se $f(n) = m$, então $\Th{Q} \Proves
\lforall[y][(!A_f(\num{n}, y) \lif \eq[y][\num{m}])]$. Mais uma vez,
esboçamos o argumento informalmente, deixando a formalização para o
leitor.

Suponha $!A_f(\num{n}, b)$. Daí obtemos (a) $!A_g(b, \num{n},
\Obj 0)$ e (b) $\lforall[w][(w < b \lif \lnot !A_g(w, \num{n}, \Obj
  0))]$. Por \olref{lem:trichotomy}, $(b < \num{m} \lor \num{m} < b)
\lor \eq[b][\num{m}]$. Mostraremos que tanto $b < \num{m}$ quanto
$\num{m} < b$ levam a uma contradição.

Se $\num{m} < b$, então obtemos
$\lnot !A_g(\num{m}, \num{n}, \Obj 0)$ de~(b). Mas $m = f(n)$; logo,
$g(m, n) = 0$ e, portanto, $\Th{Q} \Proves
!A_g(\num{m}, \num{n}, \Obj 0)$, pois $!A_g$ representa~$g$. Temos,
assim, uma contradição.

Suponha agora $b < \num{m}$. Como $\Th{Q} \Proves \lforall[w][(w <
  \num{m} \lif \lnot !A_g(w, \num{n}, \Obj 0))]$ por
\olref{rep-less}, obtemos $\lnot !A_g(b, \num{n}, \Obj 0)$. Isso
contradiz novamente~(a).
\end{proof}

\end{document}
